% !TeX root = ../main.tex
\begin{center}\textbf{ภาคผนวกประกอบบทความระบบลานิสิต}
\end{center}
รายละเอียดฟังก์ชัน รายงาน การประเมิน และแผนภาพฐานข้อมูล
\appendix
\renewcommand{\thesection}{\ifcase\value{section}\or ก\or ข\or ค\or
ง\or จ\or ฉ\or ช\or ซ\fi}
\section{ฟังก์ชันและรายงานตามบทบาทผู้ใช้งาน}
นิยามบทบาทใช้ชื่อและรหัสเดียวกันทั้งเล่ม: นิสิต (STUDENT),
อาจารย์ที่ปรึกษา (ADVISOR), หัวหน้าภาควิชา (DEPARTMENT\_HEAD),
รองคณบดี (ASSOCIATE\_DEAN), อาจารย์ผู้สอน (LECTURER) และผู้ดูแลระบบ
(ADMIN; บริการรายงานยอมรับ ADMIN\_STAFF ด้วย)
บทบาทอ่านรายงานอย่างเดียวคือ DEPARTMENT\_HEAD\_REPORT และ
EXECUTIVE\_REPORT ซึ่งไม่ให้สิทธิ์อนุมัติจากการอ่านรายงาน
ผู้ใช้หนึ่งคนอาจมีหลายบทบาท
แต่ต้องผ่านการตรวจตัวตนและความสัมพันธ์กับข้อมูลทุกครั้ง
FR06 หมายถึงงานผู้ดูแลระบบใบลา ส่วน FR07
เป็นงานผู้ดูแลข้อมูลในระบบข้อมูลสนับสนุน
ซึ่งต้องได้รับสิทธิ์ในระบบนั้นแยกต่างหาก ไม่อนุมานจาก ADMIN
ของระบบใบลา เครื่องหมาย {\fontspec{Segoe UI Symbol}✓}
หมายถึงใช้ได้ภายในขอบเขตที่ได้รับสิทธิ์ และ --
หมายถึงไม่มีสิทธิ์จากบทบาทนั้นโดยตรง
\begingroup
\setlength{\tabcolsep}{2pt}
\renewcommand{\arraystretch}{1.15}
\begin{table}[H]
  \caption{ฟังก์ชันของระบบจำแนกตามบทบาท}\label{tab:appendix-functions}
  \centering
  \begin{tabularx}{\textwidth}{|X|*{6}{>{\centering\arraybackslash}p{15mm}|}}
    \hline
    \textbf{ฟังก์ชัน} & \textbf{นิสิต} & \textbf{อ.ที่ปรึกษา} &
    \shortstack{\textbf{หัวหน้า}\\\textbf{ภาควิชา}} &
    \shortstack{\textbf{อาจารย์}\\\textbf{ผู้สอน}} &
    \textbf{รองคณบดี} & \shortstack{\textbf{ผู้ดูแล}\\\textbf{ระบบ}} \\ \hline
    ยื่นคำขอลาและแนบหลักฐาน & {\fontspec{Segoe UI Symbol}✓} & -- & --
    & -- & -- & -- \\ \hline
    ติดตามสถานะและดูประวัติใบลาของตนเอง & {\fontspec{Segoe UI
    Symbol}✓} & -- & -- & -- & -- & -- \\ \hline
    แก้ไขหลักฐานตามสถานะที่ระบบอนุญาต & {\fontspec{Segoe UI Symbol}✓}
    & -- & -- & -- & -- & -- \\ \hline
    เปิดหรือดาวน์โหลดเอกสารใบลาที่มีสิทธิ์เข้าถึง & {\fontspec{Segoe
    UI Symbol}✓} & {\fontspec{Segoe UI Symbol}✓} & {\fontspec{Segoe
    UI Symbol}✓} & {\fontspec{Segoe UI Symbol}✓} & {\fontspec{Segoe
    UI Symbol}✓} & {\fontspec{Segoe UI Symbol}✓} \\ \hline
    พิจารณาอนุมัติหรือปฏิเสธใบลาขั้นปัจจุบัน & -- & {\fontspec{Segoe
    UI Symbol}✓} & {\fontspec{Segoe UI Symbol}✓} & -- &
    {\fontspec{Segoe UI Symbol}✓} & -- \\ \hline
    ตรวจสอบประวัติการพิจารณาของตนเอง & -- & {\fontspec{Segoe UI
    Symbol}✓} & {\fontspec{Segoe UI Symbol}✓} & -- & {\fontspec{Segoe
    UI Symbol}✓} & -- \\ \hline
    ดูข้อมูลนิสิตในความดูแลของที่ปรึกษา & -- & {\fontspec{Segoe UI
    Symbol}✓} & -- & -- & -- & -- \\ \hline
    ตั้งค่ากฎอนุมัติอัตโนมัติในขั้นที่รับผิดชอบ & -- &
    {\fontspec{Segoe UI Symbol}✓} & {\fontspec{Segoe UI Symbol}✓} &
    -- & {\fontspec{Segoe UI Symbol}✓} & -- \\ \hline
    ดูรายการใบลาในขอบเขตหน่วยงาน & -- & -- & {\fontspec{Segoe UI
    Symbol}✓} & -- & {\fontspec{Segoe UI Symbol}✓} & {\fontspec{Segoe
    UI Symbol}✓} \\ \hline
    จัดการปีการศึกษาและภาคการศึกษา & -- & -- & -- & -- & -- &
    {\fontspec{Segoe UI Symbol}✓} \\ \hline
    จัดการข้อมูลคณะและสาขา & -- & -- & -- & -- & -- &
    {\fontspec{Segoe UI Symbol}✓} \\ \hline
    รับทราบใบลาที่อนุมัติแล้วของรายวิชาที่สอน & -- & -- & -- &
    {\fontspec{Segoe UI Symbol}✓} & -- & -- \\ \hline
  \end{tabularx}
\end{table}
การพิจารณาทำได้เฉพาะผู้รับผิดชอบขั้นปัจจุบันและตามลำดับที่กำหนด
การแก้ไขหลักฐานขึ้นกับสถานะคำขอ งานข้อมูลสนับสนุนตาม FR07
ต้องมีสิทธิ์แยกต่างหาก; ผู้สอนไม่มีสิทธิ์อนุมัติจากบทบาท LECTURER
เพียงอย่างเดียว
\begin{table}[H]
  \caption{รายงานของระบบจำแนกตามบทบาท}\label{tab:appendix-reports}
  \centering
  \begin{tabularx}{\textwidth}{|X|*{6}{>{\centering\arraybackslash}p{15mm}|}}
    \hline
    \textbf{รายงาน} & \textbf{นิสิต} & \textbf{อ.ที่ปรึกษา} &
    \shortstack{\textbf{หัวหน้า}\\\textbf{ภาควิชา}} &
    \shortstack{\textbf{อาจารย์}\\\textbf{ผู้สอน}} &
    \textbf{รองคณบดี} & \shortstack{\textbf{ผู้ดูแล}\\\textbf{ระบบ}} \\ \hline
    ประวัติและสถิติการลาของตนเอง แยกเทอมและรายวิชา & {\fontspec{Segoe
    UI Symbol}✓} & -- & -- & -- & -- & -- \\ \hline
    ภาพรวมและรายละเอียดการลาของนิสิตในที่ปรึกษา & -- &
    {\fontspec{Segoe UI Symbol}✓} & -- & -- & -- & -- \\ \hline
    ภาพรวมระดับภาควิชา สรุปตามอาจารย์ นิสิต และรายวิชา & -- & -- &
    {\fontspec{Segoe UI Symbol}✓} & -- & -- & -- \\ \hline
    ผลกระทบรายวิชาและเวลาเฉลี่ยการพิจารณาระดับภาควิชา & -- & -- &
    {\fontspec{Segoe UI Symbol}✓} & -- & -- & -- \\ \hline
    ภาพรวมระดับคณะ เปรียบเทียบภาควิชาและหลักสูตร & -- & -- & -- & --
    & {\fontspec{Segoe UI Symbol}✓} & -- \\ \hline
    สรุปตามอาจารย์และเวลาเฉลี่ยการพิจารณาระดับคณะ & -- & -- & -- & --
    & {\fontspec{Segoe UI Symbol}✓} & -- \\ \hline
    ภาพรวมและแนวโน้มระดับภาควิชาสำหรับผู้ดูแลระบบ & -- & -- & -- & --
    & -- & {\fontspec{Segoe UI Symbol}✓} \\ \hline
    สถิติและรายละเอียดการลารายวิชาในความรับผิดชอบ & -- & -- & -- &
    {\fontspec{Segoe UI Symbol}✓} & -- & -- \\ \hline
  \end{tabularx}
\end{table}
\endgroup
รายงานทุกส่วนจำกัดข้อมูลตามตัวตนและความสัมพันธ์กับหน่วยงาน
รายงานผู้ดูแลระบบปัจจุบันอยู่ในภาควิชาตามหน่วยงานของบัญชี
ไม่ครอบคลุมทั้งคณะ บทบาท DEPARTMENT\_HEAD\_REPORT และ
EXECUTIVE\_REPORT
อาจได้รับสิทธิ์อ่านรายงานเฉพาะด้านโดยไม่มีสิทธิ์พิจารณาใบลา
ข้อมูลรายงานอ้างอิงรอบโหลดคลังข้อมูลที่สำเร็จล่าสุด
\section{เกณฑ์ยอมรับของข้อกำหนดการทำงาน}
เกณฑ์ต่อไปนี้เป็นผลคาดหวังสำหรับทดสอบ ยังไม่ใช่ผลผ่านจริง
ควรทดสอบทั้งกรณีอนุญาตและปฏิเสธสิทธิ์โดยเปลี่ยนตัวตนผู้ใช้จริง
\begin{tabularx}{\textwidth}{|p{16mm}|X|}
  \hline
  รหัส & เกณฑ์ยอมรับ \\ \hline
  FR01 & ข้อมูลครบสร้างคำขอ PENDING และขั้นแรก CURRENT;
  ข้อมูลจำเป็นไม่ครบต้องถูกปฏิเสธโดยไม่เกิดคำขอบางส่วน \\ \hline
  FR02 & เจ้าของอ่านใบลาตนเองได้ แต่เข้าถึงของผู้อื่นไม่ได้;
  แก้ไขหลักฐานได้เฉพาะสถานะที่อนุญาต; เอกสารตรงกับคำขอ \\ \hline
  FR03 & ผู้รับผิดชอบขั้นปัจจุบันตัดสินได้;
  ผู้ที่ไม่เกี่ยวข้องหรือยังไม่ถึงขั้นตัดสินไม่ได้;
  รายงานจำกัดนิสิตในความดูแล \\ \hline
  FR04 & อนุมัติระหว่างทางเปิดขั้นถัดไป; ขั้นสุดท้ายเปลี่ยนเป็น
  APPROVED; ปฏิเสธเปลี่ยนเป็น REJECTED; รายงานคงขอบเขตหน่วยงาน \\ \hline
  FR05 & รับทราบเฉพาะใบลาที่อนุมัติและรายวิชาที่มีความสัมพันธ์;
  การรับทราบไม่เปลี่ยนผลอนุมัติและไม่เพิ่มยอดใบลา \\ \hline
  FR06 & การอ่านรายละเอียดและจัดการภาคการศึกษา
  คณะและสาขาต้องตรงกับสิทธิ์ผู้ดูแลระบบ;
  เปลี่ยนค่าบทบาทหรือรหัสผู้ใช้ในคำขอแล้วต้องไม่ขยายสิทธิ์ \\ \hline
  FR07 & นำเข้าแบบเพิ่มไม่สร้างรายการซ้ำตามคีย์;
  แบบแทนที่ได้ชุดข้อมูลตามแฟ้ม; รหัสและความสัมพันธ์ถูกต้อง;
  แถวผิดรูปแบบต้องตรวจพบ \\ \hline
  FR08 & ส่งผลซ้ำและตัดสินพร้อมกันได้ผลเดียว; ตัดสินผิดลำดับไม่ได้;
  แยกแหล่ง MANUAL และกฎอัตโนมัติย้อนหลังได้ \\ \hline
  FR09 & โหลดซ้ำไม่เพิ่มยอด; ล้มเหลวไม่เผยแพร่ batch บางส่วน;
  หนึ่งใบหลายวิชายังนับหนึ่งใบ; เทอมไม่มีใบลามีค่าเป็นศูนย์ \\ \hline
\end{tabularx}
\subsection{การบันทึกหลักฐานและการทดสอบซ้ำ}
หลักฐานของแต่ละกรณีควรประกอบด้วยรุ่นซอร์สโค้ด สภาพแวดล้อม บทบาทผู้ใช้
ข้อมูลตั้งต้น ผลที่คาดหวัง ผลที่พบ และเวลาเก็บข้อมูล หากเกี่ยวข้องกับ
warehouse ให้ระบุ batch ที่อ่านและเวลาการโหลดสำเร็จล่าสุดด้วย
เพื่อแยกกรณีประมวลผลผิดออกจากกรณีข้อมูลรายงานยังไม่ปรับปรุง
การทดสอบสิทธิ์ควรมีทั้งกรณีที่ได้รับอนุญาตและกรณีที่ต้องถูกปฏิเสธ
ผู้ทดสอบต้องเปลี่ยนตัวตนและความสัมพันธ์กับนิสิต
ไม่เพียงเปลี่ยนชื่อบทบาทในหน้าจอ เมื่อบริการอ้างอิงไม่พร้อมใช้งาน
ควรตรวจว่าระบบตอบความผิดพลาดอย่างเหมาะสมและไม่ขยายขอบเขตข้อมูลโดยอาศัยค่าที่ผู้เรียกส่งเอง
ภาพหน้าจอใช้ยืนยันสิ่งที่ผู้ใช้เห็นได้
แต่ไม่เพียงพอสำหรับยืนยันธุรกรรมหรือการป้องกันการตัดสินซ้ำ
จึงควรเก็บผลเรียกบริการและผลตรวจข้อมูลที่จำเป็นควบคู่กัน
โดยปิดบังรหัสส่วนบุคคลและข้อมูลลับ
การจัดเก็บหลักฐานในรูปแบบเดียวกันช่วยให้ตรวจซ้ำเมื่อแก้กฎหรือเปลี่ยนข้อมูลอ้างอิงได้
\subsection{แผนประเมินประโยชน์และข้อจำกัด}
การประเมินกับผู้ใช้ควรวัดเวลาที่ใช้ยื่นคำขอ ความสำเร็จในการทำงาน
และความพึงพอใจ แยกกลุ่มนิสิต อาจารย์ และผู้ดูแลระบบ
พร้อมระบุจำนวนผู้เข้าร่วมและเงื่อนไขเปรียบเทียบ
หากประเมินบริการควรรายงานเวลาโต้ตอบ เปอร์เซ็นไทล์
และอัตราข้อผิดพลาดภายใต้ภาระงานที่ระบุ
บทความฉบับนี้ยังไม่มีข้อมูลดังกล่าวจึงไม่รายงานค่าประสิทธิภาพ
และไม่ได้อ่านแถวธุรกรรมจริง จึงยังไม่ทราบจำนวนใบลา จำนวนผู้ใช้
หรือสัดส่วนประวัติที่ขาด
การมีคอลัมน์เวลาในสคีมาไม่ยืนยันว่าข้อมูลย้อนหลังครบ
ต้องเก็บสถิติรวมที่ปิดบังข้อมูลส่วนบุคคลจากสภาพแวดล้อมใช้งานเพิ่มเติม
ข้อจำกัดที่เอกสารโครงการระบุคือข้อมูลเจ้าของรายวิชา ภาควิชา
และประเภทหลักสูตรยังต้องมี mapping ทางการเพื่อให้รายงานครบถ้วน
ค่าที่ไม่ทราบคงไว้เป็นข้อมูลไม่ทราบ ไม่แทนด้วยสัดส่วนจำลอง
ระบบรองรับการบันทึกเวลาเริ่มพิจารณาและเวลาตัดสินแล้ว
แต่ยังต้องตรวจความครบถ้วนของประวัติจริง
จึงไม่ควรแทนระยะเวลาที่ขาดด้วยศูนย์
ระบบข้อมูลสนับสนุนให้ข้อมูลอ้างอิงปัจจุบัน
ขณะที่ระบบใบลาเก็บตารางเรียนตามภาคการศึกษา
จึงต้องระวังการผูกข้อมูลนี้กับประวัติย้อนหลัง
อีกทั้งยังต้องกำหนดรอบโหลดข้อมูลรายงานอัตโนมัติให้เหมาะกับสภาพแวดล้อมใช้งาน

\subsection{เส้นทางการพิจารณาและการสิ้นสุดคำขอ}
\begin{figure}[H]\centering
  \begin{tikzpicture}[box/.style={draw,align=center,text
        width=68mm,minimum
    height=12mm,font=\fontsize{12bp}{16bp}\selectfont},flow/.style={-{Stealth},thick}]
    \node[box] (advisor) at (0,0) {ที่ปรึกษา CURRENT\\ใบลา PENDING};
    \node[box] (head) at (0,-2.1) {หัวหน้าภาควิชา CURRENT\\ใบลา PENDING};
    \node[box] (dean) at (0,-4.2) {รองคณบดี CURRENT\\เฉพาะวันลามากกว่า 15 วัน};
    \node[box] (done) at (0,-6.3) {อนุมัติครบทุกขั้น: ใบลา APPROVED};
    \draw[flow] (advisor)--node[right]{อนุมัติ}(head);
    \draw[flow] (head)--node[right]{อนุมัติและเกิน 15 วัน}(dean);
    \draw[flow] (dean)--node[right]{อนุมัติ}(done);
    \draw[flow]
    (head.west)--(-4.4,-2.1)--node[left,align=right]{อนุมัติ\\ไม่เกิน
    15 วัน}(-4.4,-6.3)--(done.west);
  \end{tikzpicture}
  \caption{เส้นทางอนุมัติตามโค้ดและเงื่อนไขจำนวนวันลา}\label{fig:approval-route}
\end{figure}
ภาพแสดงลำดับพิจารณาเมื่อข้อมูลผู้รับผิดชอบครบ
จึงควรตรวจสอบที่ปรึกษาและผู้รับผิดชอบแต่ละระดับก่อนเริ่มใช้งาน
ทุกขั้นปฏิเสธได้เฉพาะเมื่อเป็น CURRENT แล้วใบลาเปลี่ยนเป็น REJECTED;
ขั้นก่อนหน้าที่อนุมัติแล้วไม่ทำให้ใบลาทั้งใบสำเร็จ
กฎอัตโนมัติใช้เส้นทางเดียวกันและแยกแหล่งคำตัดสิน ผู้สอนรับทราบภายหลัง
APPROVED โดยไม่เปลี่ยนผลอนุมัติ
% !TeX root = ../main.tex
\clearpage
\begin{landscape}
  \setlength{\intextsep}{4bp}
  \section{แผนภาพฐานข้อมูลระบบใบลา}\label{app:db}
  \begin{figure}[H]
    \centering
    \includegraphics[width=\linewidth,height=148mm,keepaspectratio]{images/leave-schema/DB.png}
    \caption{โครงสร้างฐานข้อมูลระบบใบลา}\label{fig:db-original}
  \end{figure}
  \clearpage
  \section{ภาพรวมคลังข้อมูลรายงาน}\label{app:dw-overview}
  \begin{figure}[H]
    \centering
    \includegraphics[width=\linewidth,height=148mm,keepaspectratio]{images/leave-schema/DW-OVERVIEW.png}
    \caption{ภาพรวมองค์ประกอบคลังข้อมูลรายงาน}\label{fig:dw-overview-original}
  \end{figure}
  \clearpage
  \section{แผนภาพความสัมพันธ์คลังข้อมูล}\label{app:dw-detail}
  \begin{figure}[H]
    \centering
    \includegraphics[width=\linewidth,height=148mm,keepaspectratio]{images/leave-schema/DW-ERD.png}
    \caption{ความสัมพันธ์ของข้อมูลในคลังข้อมูลรายงาน}\label{fig:dw-erd-original}
  \end{figure}
  \clearpage
  \section{แผนภาพคลังข้อมูลฉบับเต็ม}\label{app:dw-full}
  \begin{figure}[H]
    \centering
    \includegraphics[width=\linewidth,height=148mm,keepaspectratio]{images/leave-schema/DW-FULL_VERSION_2.png}
    \caption{โครงสร้างคลังข้อมูลรายงานฉบับเต็ม}\label{fig:dw-full-original}
  \end{figure}
  \clearpage
  \section{แผนภาพส่วนควบคุมคลังข้อมูล}\label{app:controls}
  \begin{figure}[H]
    \centering
    \includegraphics[width=\linewidth,height=148mm,keepaspectratio]{images/leave-schema/DW-CONTROL_COLUMNS.png}
    \caption{ส่วนควบคุมการปรับปรุงข้อมูลในคลังข้อมูล}\label{fig:dw-control-original}
  \end{figure}
\end{landscape}


\clearpage
\section{ตัวอย่างหน้าจอระบบลานิสิต}\label{app:screenshots}
\begin{figure}[H]
  \centering
  \includegraphics[width=\linewidth,height=91mm,keepaspectratio]{work_capture/images/lecturer-request-detail-legacy.png}
  \caption{หน้ารายละเอียดคำร้องสำหรับอาจารย์ผู้สอน
  แสดงผลการพิจารณาและรายวิชาที่ได้รับผลกระทบ}\label{fig:screen-lecturer-detail}
\end{figure}
\vspace{8mm}
\begin{figure}[H]
  \centering
  \includegraphics[width=\linewidth,height=91mm,keepaspectratio]{work_capture/images/student-document.png}
  \caption{หน้าแสดงเอกสารคำร้องขอลาและปุ่มดาวน์โหลด
  PDF}\label{fig:screen-student-document}
\end{figure}
\clearpage
\noindent\textbf{ตัวอย่างหน้าจอระบบลานิสิต (ต่อ)}
\begin{figure}[H]
  \centering
  \includegraphics[width=\linewidth,height=91mm,keepaspectratio]{work_capture/images/student-report-2-legacy.png}
  \caption{หน้ารายงานของนิสิตแสดงจำนวนการลาแยกตามรายวิชา}\label{fig:screen-student-report}
\end{figure}
\vspace{8mm}
\begin{figure}[H]
  \centering
  \includegraphics[width=\linewidth,height=91mm,keepaspectratio]{work_capture/images/department-report-advisor-processing.png}
  \caption{หน้ารายงานเวลาในการพิจารณาขั้นตอนที่ 1
  แยกตามอาจารย์ในภาควิชา}\label{fig:screen-advisor-processing}
\end{figure}
